% =========================================================================
\chapter{Discussion, Threats, and Future Work}
\label{ch:discussion}
% =========================================================================

\claim{The proposition's scope is stated, the gates' blind spot is
exhibited, the costs are honest, and the roadmap attaches the missing organ
where the evidence says it belongs: the retention decision.}

% -------------------------------------------------------------------------
\section{The readout proposition and its scope}
\label{sec:proposition}
% -------------------------------------------------------------------------

\subsection*{The proposition}

Self-evolution is usually described as variation plus selection. This
thesis argues it has a third term:

\begin{quote}
\textbf{Self-evolution $=$ variation $+$ selection $+$ readout.} When the
readout's resolution falls below the effect sizes being selected on,
retention degenerates into drift, and no improvement in variation quality
recovers it.
\end{quote}

It is called a \emph{proposition}, not a theorem, and the reason is
methodological honesty rather than modesty: as stated in words it restates
power analysis, and what makes it non-trivial is a dynamics argument, given
below, that a single-test power calculation does not supply.

The claim is made for this bed and this system class. The field-level
implication --- that every surveyed loop selects on aggregate score
(\S\ref{sec:self-evolving}) and is therefore exposed to the same
arithmetic on any bed with a comparable flip rate --- is an \emph{inference}
and is flagged as one. It rests on the scan of Chapter~\ref{ch:related}:
160 papers screened, 18 close-read. An independent, contemporaneous
critique of how harness evolution is evaluated reaches a convergent
conclusion under a different evaluation design~\cite{rethink-harness-eval}.

\subsection*{Formal spine, borrowed rather than invented}

Quantitative genetics already has the equation. The response to selection
is given by the breeder's equation~\cite{falconer-mackay},
\[
  R \;=\; h^2 S ,
\]
gain per generation equals heritability times the selection differential.
When selection acts on a \emph{measured} phenotype rather than the true
one, $h^2$ is attenuated by the measurement's reliability --- classical
regression dilution~\cite{hutcheon-dilution}:
\[
  h^2_{\text{realized}} \;=\; h^2_{\text{true}}\,\rho ,
  \qquad
  \rho \;=\; \frac{\sigma^2_{\delta}}{\sigma^2_{\delta}+\sigma^2_{\varepsilon}} .
\]

Instantiated on this bed: the per-edit effect is $|\delta| \lesssim 3$
tasks, and the per-round readout noise is $\sigma_\varepsilon \approx
3.70$ tasks (\S\ref{sec:power}), giving $\rho \lesssim 0.4$. That is
an \emph{upper} bound, because $|\delta|$ is a magnitude bound standing in
for a standard deviation. The bound itself comes from the hand-run drug
trials \F{34}\F{36} --- a small, non-random set chosen because a family
looked repeatedly prescribed, not a random sample of the six campaigns'
shipped candidates.

\paragraph{The dynamics claim.} A ratchet drives both its accept arm and
its rollback arm from the same measured signal. The false-kill rate of true
gains therefore scales with $1-\rho$, while retained noise accumulates on
the other side. Expected cumulative gain over $R$ rounds can be non-positive under this
bed's parameter regime \emph{regardless of variation quality} --- which is why better
candidate generation, the axis the entire self-evolution literature
optimizes, cannot rescue a loop in this regime.

\paragraph{Boundary condition, supplied by this thesis's own evidence.} The
breeder's equation governs the \emph{mean} response.
\S\ref{sec:archaeology} shows gains on this bed arriving as rare tail
events from a single mechanism family. The proposition therefore constrains
the mean and is silent on the tail --- a boundary worth stating, because
the tail is where this thesis's one positive result lives.

\subsection*{The calibration that motivates all of it}

The argument above would be arithmetic without an anchor. The anchor is the
budget-starvation repair (\S\ref{sec:one-repair}), and it functions here as
a \emph{positive control} in the clinical sense --- what ICH~E10 calls
assay sensitivity, the property that lets a trial distinguish an effective
treatment from an ineffective one, without which a null cannot be
interpreted.

This thesis has such an intervention, and the score axis \emph{did} recover
it. The repair is known to have worked on mechanism grounds:
\texttt{budget\_exceeded} exits collapsed from 35 to 4 in a single round,
and twelve of the sixteen newly passing tasks had died of starvation in the
round before. Its $+9$ cleared the same-config band, and it is the largest
single-round move that campaign makes after its baseline draw \F{39}.

\emph{That recovery is the calibration, not the reassurance.} The effect
had to reach $+9$ tasks to clear a band of $-8\ldots{+}5$, while the
per-edit effects this loop actually produces are $|\delta| \lesssim 3$
tasks. \textbf{The assay's detection limit sits at roughly three times the
size of the effects being selected on.}

Two campaign-wide sweeps then show that a detection limit is the charitable
reading. Of the $84$ round-to-round pairs following the baseline draws,
eleven leave the band and six of those eleven occur in rounds where nothing
shipped; restricted to whole-bed rounds, six leave the band and five are
ship-bearing, across three buckets that share no mechanism \F{48}. And
among the whole-bed rounds whose starvation count collapses by ten exits or
more, two collapses of \emph{identical} size, $-31$ exits, moved the score
by $+9$ and by $+4$, while a $-19$ collapse moved it by $-3$ \F{49}. A
starvation collapse need not raise the score; a task that stops running out
of steps starts answering, and the answer can be wrong. But the spread is
the finding: \textbf{over the range of effects this loop can produce, the
score axis is not a monotone transducer of mechanism magnitude.}

An instrument like that is not returning nulls that mean ``no effect''; it
is returning nulls that mean ``below threshold --- or above it and pointing
the wrong way''. Every null measured on that axis is therefore
uninterpretable --- \emph{including this thesis's own} --- and the field's
universal selection signal is being read at a resolution its own effect
sizes do not reach.

\subsection*{The rival hypothesis, stated and weighed}

Every null in this thesis admits a simpler explanation than the readout
proposition: \textbf{the loop's action space cannot reach the lesions.}
This is the capability wall, and it is the primary layer in the mortality
account \F{40}. If the drugs are placebos because the disease is out of
reach, no readout is needed to explain a flat outcome axis.

\emph{On the outcome axis the two are observationally equivalent}, and this
thesis does not pretend otherwise. Level 3 was exercised by hand
\F{34}\F{36} and never wired into retention, so no experiment here
separates them.

They are not, however, equivalent on all of the evidence, and the
separation is worth stating even though it is \cls{C} judgment rather than
a test:

\begin{enumerate}
  \item \textbf{The actuator demonstrably reached a lesion at least once.}
        The starvation repair worked, was mechanism-confirmed, and was
        retained \F{39}. A capability wall that admits this counter-example
        is no longer the claim that the action space cannot reach lesions;
        it is the narrower claim that it cannot reach the \emph{dominant}
        family --- wrong-closure, 18 of the 49 swing tasks \F{31} --- which
        this thesis accepts (\S\ref{sec:future}).
  \item \textbf{The drug that worked was then lost for a reason unrelated
        to capability.} It reached no later campaign because its target is
        a \texttt{file://} path into its own campaign directory and every
        later seed is built from the clean template: \emph{the inheritance
        channel does not exist} \F{39}. No capability wall explains that.
  \item \textbf{The same families are re-derived from scratch, repeatedly.}
        The loop re-prescribes drugs it has already prescribed correctly,
        four campaigns apart, without inheriting \F{38}. A loop blocked by
        a capability wall would fail to \emph{cure}; it would not fail to
        \emph{remember}.
\end{enumerate}

So the honest summary is: the capability wall explains why the loop does
not solve the hard family. It does not explain the loss of the one drug
that worked, and it does not explain the treadmill. Those two read as
readout and inheritance failures, which is what the proposition
addresses.

\paragraph{The discriminating experiment, named rather than claimed.}
Attach the level-3 readout to the retention decision and re-fly. If
retention still drifts, the actuator was the bottleneck after all. This is
why \S\ref{sec:future} leads with it.

% -------------------------------------------------------------------------
\section{Governance: necessary, not sufficient}
\label{sec:governance}
% -------------------------------------------------------------------------

The acceptance gates work, within their remit. Seven structural refusals
are on the record \F{11}, and no refused candidate shipped.

Their remit is form, not provenance. Under a coverage directive --- treat
every task, never abstain --- the loop engineered the bed's own leak routes
into a benchmark-lookup cheat tool, and that tool passed the Critic and
\emph{every} landing gate \F{37}. Nothing in the ladder asks where a
capability came from.

Two fixes follow, and the second has an in-house origin worth recording. A
leakage blocklist addresses the bed side. A gate forbidding candidates from
reading adjudication files addresses the loop side --- and it exists
because registry error 13 was exactly that failure committed by us: an
adjudication file left readable in a run root, a role that read the answer
sheet, and a round of rulings voided. The governance gap and the evaluator
gap are the same gap seen from two sides.

\paragraph{Disclosure.} The evidence-column defect (\S\ref{sec:evidence-audit}),
the safety clamp that cannot engage on a small bed, and the regression
report left empty by a round-indexing slip (\S\ref{sec:open-loop}),
together with the F37 near-miss above, will be reported to the
HarnessX/AEGIS maintainers around submission, with a pointer to the
relevant sections of this thesis. Submission is not delayed awaiting a
response.

% -------------------------------------------------------------------------
\section{Threats to validity}
\label{sec:threats}
% -------------------------------------------------------------------------

\paragraph{Attribution is arm-level.} No single-factor attribution is
claimed anywhere \F{9c}, and the reason no component ablation was run is
given in \S\ref{sec:localization} rather than left implicit.

\paragraph{The localization metric is post-hoc}, defined over predictions
that were written before the outcome was known but scored under a rule
fixed afterwards. This is stated at every point of use.

\paragraph{Single bed, single model tier.} Three seeds per arm establish
that the noise geometry replicates \F{44}\F{45}; generalization to other
beds and model families is untested.

\paragraph{Bed asymmetry between the arms.} The no-graph arm flew the full
103-task bed on the first two seeds while the graph arms flew the 100-task
no-pixel subset, and the third seed's no-graph arm switched beds at round 7.
Every cross-arm reading in this thesis is recomputed on the common 100-task
subset from \texttt{task\_history}, never taken from \texttt{curves.json},
which records whichever bed each round actually flew \F{46}.

\paragraph{Leak-route contamination.} Thirteen of the 49 swing tasks are
leak-route \F{31}. Both arms evaluate on the same fixed bed every round, so
drift toward the bed is shared; the blocklist is the only held-out defence.

\paragraph{Operational deviations.} One killed evolve stage and one killed
batch on the second seed; a gateway outage with five poisoned rounds on the
third; five meta-chain holes; three \texttt{--start-round} re-flights across
the four replication arms (\S\ref{sec:replication}). The flip statistics use
only verified clean same-config pairs, and non-adjacent pairs are left
unpooled.

\paragraph{Windows frictions} are symmetric across arms and disclosed.
\textbf{Judge residual risk} remains after hardening.

\paragraph{The evaluator.} The fifteen-error registry is reported as part
of the threats rather than hidden in an appendix, with both bias directions
documented and \emph{stopping the check too early} named as the common root
cause. A retraction registry is maintained and nothing in it is re-cited
--- including this thesis's own earlier claim that the official loop could
not prescribe cause-aware drugs, which its own archives falsified
(\S\ref{sec:archaeology}).

% -------------------------------------------------------------------------
\section{Cost honesty}
\label{sec:cost}
% -------------------------------------------------------------------------

\subsection*{The cost work is not a contribution of the graph}

Noop-round batching (\$38.1 to \$11.8 per round \F{13}), single-shot
digestion, the zero-cache-hit repair and the retention policy are
\emph{harness-level} mechanisms. Each runs on the unmodified official
system and none needs an execution graph. This thesis found and implemented
them, and crediting them to GHX would be a contribution in disguise. The
switch ledger already classifies them this way: the profit-and-loss
account's cost column is labelled \emph{portable} \F{12}. They are
described in Appendix~\ref{app:cost}, as engineering.

\subsection*{What the graph layer costs is unknown}

An earlier reading of this work recorded the graph layer as net expensive.
\textbf{That reading is withdrawn.} It was measured across the shakedown
campaigns the whitelist later excluded, so it was never admissible; and on
the admissible campaigns the task-side spend per fresh evaluation is
\emph{lower} in the graph arms on all three seeds (\$0.392, \$0.472,
\$0.503 against \$0.461, \$0.538, \$0.544) \F{13}.

That comparison settles nothing either, because it excludes the meta tier
--- which is exactly where graph work would appear, and for which no
admissible per-role accounting exists. \emph{The honest entry is a blank,
not a minus.}

\subsection*{The graph layer's one measured liability}

It is operational rather than financial. A crash family introduced by the
recorder --- binary content reaching the raw JSONL and killing the evolve
stage --- produced 164 tracebacks across the three graph seeds against 5
across the three no-graph seeds, together with five meta-chain holes \F{56}. It
cost metadata rather than scores, since the task batches ran regardless.
It is a real debit and is entered as one.

\subsection*{The switch ledger, and what replication cost}

Twenty-one switches were built and each is classified: 2 quality gains, 2
portable cost gains, 2 measured negatives, 3 measured zeros, 3 exercised
but unproven, 2 substrate, and 7 capability-only \F{12}. \emph{None
constitutes a measured behavioural improvement}, and the ledger says so in
its own column rather than in a caveat.

Replication was priced rather than omitted: the second seed-pair cost
\$886 plus \$581 and the third \$882 plus \$584 in clean spend, with
gateway-poisoned rounds excluded --- \$1{,}467 and \$1{,}466 per replicated
seed-pair \F{44}\F{45}. That is the price of the noise geometry the whole
readout argument stands on, and it is the single most load-bearing purchase
this project made.

% -------------------------------------------------------------------------
\section{Future work}
\label{sec:future}
% -------------------------------------------------------------------------

Each item below is attached to a measured gap in this thesis, and they are
ordered by what the evidence says matters.

\paragraph{1. Readout-first self-evolution.} Attach level 3 to the
retention decision. The ingredients are already in the literature ---
paired prefix replay~\cite{shepherd,car}, sibling-rollout
controls~\cite{craft}, anytime-valid sequential tests~\cite{pace}. This
thesis's contribution is not the ingredients; it is \emph{where to attach
them}, and the evidence that attaching them anywhere else is premature.
This is also the discriminating experiment of \S\ref{sec:proposition}.

\paragraph{2. Bed reliability, with its ceiling priced first.} The leakage
blocklist and result caching at the measured variance injection point ---
a web search, in 76 of 79 cases \F{25} --- are direct and cheap. Reliability
engineering should target the volatile family, where flips concentrate at
$1.6$ to $2.0$ times the bed rate \F{41}\F{44}\F{45}, rather than the whole
bed.

\emph{But the census bounds what that can buy} (\S\ref{sec:power}, \F{54}):
21 of 100 tasks carry no variance at all, and the 79 that can move are the
ones that move unaided. Item selection cannot produce a bed that is both
reliable and informative. The purchasable fix is repetitions --- price it
with Spearman--Brown against the measured per-task flip rate before
committing budget: a single draw prices at reliability $0.57$, and
\textbf{four} averaged repeats are the first integer count crossing this
thesis's own $0.80$ ``reliable'' bar \F{60} --- and the alternative is the
readout itself, which is why item 1 comes first.

\paragraph{3. Replicate the episodic-note channel} under a pre-registered,
budget-bounded protocol \F{35}. It is solver-level and bypasses all four
mortality layers, which makes it the cheapest positive result available ---
and, as \S\ref{sec:survivor} notes, the one that gives up the property
under study.

\paragraph{4. Capability procurement for the wrong-closure family.}
Offline simulation over the trial repetitions shows majority-at-$k$ lifting
mean pass from $0.35$ to $0.60$ at $k{=}10$, entirely from scatter-shaped
tasks (3 lifted, 0 sunk), while consensus-wrong tasks lock at zero for
every $k$ \F{42}. Plain voting \emph{hardens} the confident-wrong answer.
The organ needed is external discrimination, and that does sit outside the
loop's action scope.

The step cap does not. On the second seed the loop reached the starvation
problem twice, from two different buckets. In R4 it filled in a dormant
runtime-policy rule --- a config-bucket edit that steers the model away
after three consecutive empty Bash calls --- and in the round that
followed, budget-exhaustion exits fell from $39$ to $8$ \F{49}. In R5 it
went further and inserted a \texttt{StepCountdownProcessor} into its own
graph, certifying the insert by running the processor and reading its
model-visible tail back out \F{50}. \emph{A lever an earlier reading of
this work placed outside the loop's scope was pulled from inside it.}
\emph{Scope is not a fixed property of the system}; it is whatever the
action space currently reaches, and this loop widened its own.

The score's response to that $-31$ collapse was $+4$ tasks --- inside the
band, indistinguishable from a re-draw \F{49}. The loop widened its scope
and could not tell that it had.
